\begin{Answer}{9.24}
$abed$ is a cycle of length 4 and $ecdab$ is a cycle of length 5.
Other answers are possible.
There is no cycle of length 6 since there are only 5 vertices in the graph and a cycle cannot
repeat a vertex.
\end{Answer}
\begin{Answer}{9.27}
You should have drawn something like this (but probably bigger and with dots on the corners): $\square$
$\triangle$
\end{Answer}
\begin{Answer}{9.30}
You should have drawn something like this (but probably bigger and with dots on the corners
and center): $\times$
\end{Answer}
\begin{Answer}{9.31}
You should have drawn a path of length 2 and 4 vertices not connected to anything. Something like this:
 $|$ . . . .
\end{Answer}
\begin{Answer}{9.47}
(a) Yes. Put the vertices on opposite corners in the same
 set.
(b) No. Every other vertex needs to be in the opposite
set, and since there is an odd cycle, eventually
we get to a situation where a vertex can't be in either set.
(c) No. None of the vertices can be in a set with
any others since they are all connected to each other.
(d) Yes. If I put $000, 110, 011, 101$ in one set and
the others in the other set, you can easily
see that all of the edges go between sets.
(e) Yes. Since it is just a path, put every other vertex in
the opposite set. Unlike the situation with the odd cycle,
we don't ever loop back so no problems can arise.
\end{Answer}
\begin{Answer}{9.48}
It is true. Notice that a tree with 2 vertices is clearly
bipartite since it is just an edge.
Assume all trees with $n$ vertices are bipartite,
where $n\geq 2$.
Let $T$ be a tree with $n+1$ vertices and $v$ be a leaf
connected to some vertex $u$.
Let $T^{\prime}$ be the tree $T$ with the leaf $v$ removed. Then by in the inductive hypothesis, $T^{\prime}$ is bipartite since it has $n$ vertices.
Then $T$ is bipartite since we can put $v$ in the
opposite set of $u$ with no risk of edges between
the nodes within that set since
 $v$ is only connected to $u$.
Thus, by PMI, all trees with $n>2$ vertices are bipartite.
\end{Answer}
\begin{Answer}{9.57}

\sputfig{graphs/fig/adjListK33}{.75}
\end{Answer}
\begin{Answer}{9.58}

\sputfig{graphs/fig/adjListK33b}{.75}
\end{Answer}
\begin{Answer}{9.63}

$
\begin{blockarray}{*{7}{c}}
 &A & B & C & D & E &E \\
\begin{block}{c[*{6}{c}]}
A&0&0&0&1&1&1  \bigstrut[t]\\
B&0&0&0&1&1&1\\
C&0&0&0&1&1&1\\
D&1&1&1&0&0&0\\
E&1&1&1&0&0&0\\
F&1&1&1&0&0&0\\
\end{block}
\end{blockarray}
$
\end{Answer}
\begin{Answer}{9.64}

$
\begin{blockarray}{*{7}{c}}
 &A & B & C & D & E &E \\
\begin{block}{c[*{6}{c}]}
A&0&0&0&1&1&1  \bigstrut[t]\\
B&0&0&0&1&1&1\\
C&0&0&0&1&1&1\\
D&0&0&0&0&0&0\\
E&0&0&0&0&0&0\\
F&0&0&0&0&0&0\\
\end{block}
\end{blockarray}
$
\end{Answer}
\begin{Answer}{9.75}
Did you draw a triangle with a vertex in the middle connected to the
three vertices of the triangle? I thought so!
\end{Answer}
\begin{Answer}{9.81}
Notice that $K_{3,3}$ does not have $C_3$ as a subgraph.
Since $K_{3,3}$ has $3\cdot 3 = 9$ edges and $9 > 8 = 2(6)-4$,
Theorem~\ref{thm:kurat} part (b) implies that $K_{3,3}$ is not planar.
\end{Answer}
\begin{Answer}{9.91}
Since $C$ is the set of edges of a connected component of
$G_A\!=\!(V,A)$,
then $A$ has no edges between vertices
of $C$ and vertices of $V-C$.

Another way to think about it: If the cut $(C,V-C)$ does not respect $A$,
then there is some $(u,v)\in A$ such that $u\in C$ and $v\in V-C$.
But $C$ is the set of edges of a {\bf\em connected} component,
and since $(u,v)\in A$, $v\in C$, contradicting the fact that $v\in V-C$.
Therefore, the cut $(C,V-C)$ respects $A$.
\end{Answer}
\begin{Answer}{9.94}
Since Idea 1 is considering adding edges adjacent to a given vertex,
Prim's algorithm is essentially using this idea, at least in the first step.
Although to be fair, Prim's algorithm is also essentially directly using
Theorem~\ref{thm:mst1}, where $C$ is the single tree it is growing.

On the other hand, Idea 2 is about using minimum weight edges
without a specific endpoint in mind, which is what Kruskal's
algorithm is doing.
\end{Answer}
\begin{Answer}{9.97}
Recall that the maximum number of edges a graph can have is
$m\leq \binom{n}{2}=\frac{n(n-1)}{2}=O(n^2)$.
Therefore, $O(\log m)=O( \log n^2)=O(2\log n)=O(\log n)$.
Given this, it is clear that $O(m\log m)=O(m\log n)$
\end{Answer}
\begin{Answer}{9.102}
Lines 1-3 and  6-7 take constant time.
Lines 4-5 line takes $O(n)$ time since we are doing $n$ things
that take constant time.
Line 8 takes $O(n)$ time.
The \verb+While+ loop executes $n$ times,
each time calling \verb+Q.extractMin+ that takes $O(\log n)$
and the \verb+ForAll+ loop.
The code in the \verb+ForAll+ loop takes $O(\log n)$
time, and it executes at most $n-1$ times (since there can be
at most $n-1$ vertices adjacent to a given vertex).
Combining this together, the complexity is
$O(n)+O(n(\log n + (n-1)\log n))=O(n^2\log n)$
(Make sure you understand how I simplified this!).

Unfortunately, this actually overestimates the time complexity
of the algorithm.
Obtaining a better bound takes a little thought and can be tricky
to understand, so read carefully!
Let's analyze the \verb+ForAll+ loop
more carefully. In particular, instead of counting up the operations
according to how the code is written (iterating over adjacency lists),
we will think about how many times the code inside that
loop executes in total regardless of when it happens.
Notice that the code in the \verb+ForAll+ loop
can execute at most two times for every edge--once from each
endpoint. If you cannot see why this is, definitely ask about it in class!
Thus, the complexity of the \verb+ForAll+ loop
is actually $O(m\log n)$ over all of the iterations of the
\verb+While+ loop. So the complexity of the \verb+While+
loop is $O(m\log n)$ plus the complexity of the \verb+Q.extractMin+ method called $n$ times, which we saw was $O(n\log n)$.
Given that, the complexity of the algorithm is
$O(n)+O(m\log n)+O(n\log n)=O(m\log n)$.
\end{Answer}
